% ---
% id: receiver/component-catalogue/rem:layers
% kind: rem
% title: Detector and source layers outside the catalogue
% summary: >-
%   Detector efficiency with dark counts, and the probe source, are declared
%   layers rather than catalogue entries. Efficiency enters as the terminal
%   factor, dark counts act on the outcome label, and source modes, unlike
%   vacuum inputs, are never eliminated.
% prominence: standard
% status: accepted
% uses: [receiver/def:compile, receiver/def:probe]
% mentions: [def:catalogue-entries, receiver/detector-imperfections/thm:efficiency, receiver/detector-imperfections/req:efficiency-schema, receiver/detector-imperfections/thm:dark-counts, receiver/thm:port-grams, receiver/rem:lossvsprobe]
% assumes: [A-VAC]
% refs: []
% provenance: newly-derived
% lean:
% ---

\begin{remark}\label{rem:layers}
Two further declared objects are \emph{layers}, not chain elements, and are
therefore not rows of Table~\ref{tab:catalogue} of
Definition~\ref{def:catalogue-entries}.

\emph{L1 --- Detector layer.} A port map $\{\Port{k}\}_{k=1}^{K}$ with
per-port efficiency $\eta_k\in[0,1]$ and dark-count probability
$p_{\dc|k}\in[0,1]$. The efficiency enters as the
\tref[terminal detector-efficiency layer]{terminal factor} $D_\eta$ of
Definition~\ref{def:compile}; the dark counts never enter $\Gmat$ at all and
act on the outcome label (Theorem~\ref{thm:dark-counts}). \emph{Detector} efficiency
is \emph{not} C7: it is the terminal layer, and the schema keeps the two apart
(Theorem~\ref{thm:efficiency}\,(iii) and Requirement~\ref{req:efficiency-schema}).

\emph{L2 --- Source (probe) layer.} The input split
$\Inp=\InpC\sqcup\InpS$ and the \tref{probe declaration} of
Definition~\ref{def:probe}. Undeclared input modes carry vacuum and may be
eliminated (Theorem~\ref{thm:port-grams}\,(iv)); \tref{declared source modes} carry a
declared state and \emph{may not}. That difference is the entire structural
distinction between ``loss'' and ``probe'', and it is a definition rather than
a convention (Remark~\ref{rem:lossvsprobe}).
\end{remark}
